% CCCG 2026 submission. Compile with pdflatex.
% Requires the cccg26.cls file from CCCG2026-tex-template/. Either run
% `cp CCCG2026-tex-template/cccg26.cls .` once, place dgf3.tex inside
% CCCG2026-tex-template/, or set TEXINPUTS to include that folder.
%
% Page limit: six pages for the main body, excluding references and
% appendix. Anything that does not fit in six pages should be in the
% appendix; the appendix is read at the discretion of the PC.

\documentclass{cccg26}
\usepackage{graphicx,amssymb,amsmath}
\usepackage{algorithm}
\usepackage{algpseudocode}
\usepackage{booktabs}

% hyperref must be loaded (nearly) last. It subsumes `url` and provides
% the \href command that bibtex emits in the .bbl for DOIs.
\usepackage[hidelinks]{hyperref}

%----------------------- Macros and Definitions --------------------------

% The cccg26 class already defines: theorem, lemma, cor, obs, conj,
% prop, invar, proof. Add only the ones we additionally need.
\newtheorem{definition}[theorem]{Definition}
\newtheorem{claim}[theorem]{Claim}

\DeclareMathOperator{\wt}{wt}
\DeclareMathOperator{\MST}{MST}
\DeclareMathOperator{\APSP}{APSP}

% Pseudocode comment style used throughout the algorithms.
\algrenewcommand{\algorithmiccomment}[1]{\hspace{3em}// #1}

%----------------------- Title -------------------------------------------

\title{Descending Greedy Filter}

% TODO: confirm affiliations and emails before submission.
\author{P.~Carmi\thanks{Department of Computer Science, Ben-Gurion University of the Negev, \texttt{carmip@cs.bgu.ac.il}}
	\and
	N.~Horovitz\thanks{Department of Computer Science, Ben-Gurion University of the Negev, \texttt{horovin@post.bgu.ac.il}}
	\and
	Y.~Yashuk\thanks{Department of Computer Science, Ben-Gurion University of the Negev, \texttt{yashuk@post.bgu.ac.il}}}

\index{Carmi, Paz}
\index{Horovitz, Noam}
\index{Yashuk, Yotam}

%------------------------------ Text -------------------------------------

\begin{document}
\thispagestyle{empty}
\maketitle


\begin{abstract}
We introduce a new approach to constructing geometric $t$-spanners based on a simple filtering principle. As part of this framework, we formalize the notion of \emph{minimal} $t$-spanners---graphs in which every edge is essential for preserving the stretch guarantee---and present the \emph{Descending Greedy Filter} (DGF), which produces such spanners.
We conjecture that all minimal $t$-spanners in the plane, for any fixed $t>1$, have $O(n)$ edges and total weight proportional to that of the minimum spanning tree.

Experimental results show that DGF produces spanners with essentially the same number of edges as the classical greedy construction, while consistently achieving $4$--$7\%$ lower total weight, without loss of stretch. To the best of our knowledge, this is the first construction that empirically outperforms the greedy spanner in total weight on random point sets.
\end{abstract}

\section{Introduction}
Geometric spanners are a central object in computational geometry. Given
a set $P$ of $n$ points in the plane and a stretch parameter $t\ge 1$, a
geometric $t$-spanner is a graph $G=(P,E)$ in which the shortest-path
distance between any pair of points is at most $t$ times their Euclidean
distance. Such structures have been extensively studied because of their
applications in network design, approximation algorithms, and distributed
systems.

Among the many known constructions, the \emph{greedy spanner} of
Alth\"ofer, Das, Dobkin, Joseph and Soares is widely regarded as the
gold standard: it achieves $O(n)$ edges, total weight
$O(\wt(\MST(P)))$, stretch exactly $t$, and, by a line of work
culminating in Filtser and Solomon, is existentially optimal in the
sense that its sparsity and weight cannot be improved in the worst
case. Empirical studies consistently place the greedy spanner ahead of
Yao graphs, $\Theta$-graphs and WSPD-based constructions on every
quality metric except construction time.

This raises a natural question: is the greedy spanner optimal not only
existentially, but also \emph{instance by instance}? Equivalently, on a
given input, can one find a $t$-spanner with lower total weight (or fewer
edges) than the greedy spanner, without sacrificing the
stretch guarantee? The greedy spanner is \emph{edge-minimal}---no
single edge can be removed while preserving the stretch---but being
edge-minimal is a much weaker property than being globally optimal, and
nothing in the greedy construction is tailored to exploit long,
redundant edges created late in the process.

\paragraph{Our contributions.}
We introduce the \emph{Descending Greedy Filter} (DGF), a construction
dual to the greedy spanner, and investigate the class of edge-minimal
$t$-spanners in the plane. Specifically:
\begin{itemize}
  \item \textbf{A new construction.} DGF (Algorithm~\ref{alg:dgf})
    starts from the complete graph and removes edges in
    \emph{non-ascending} order of length whenever doing so preserves
    the stretch $t$. The output is by construction an edge-minimal
    $t$-spanner.
  \item \textbf{A filtering procedure.} We extract from DGF a general
    filter (Algorithm~\ref{alg:dgf-filter}) that converts any
    $t$-spanner into an edge-minimal one, and prove
    (Theorem~\ref{thm:filter-minimal}) that the result is a subset of
    the input and minimal in a pair-witness sense.
  \item \textbf{Conjectures on minimality.} We state four conjectures
    (Conjectures~\ref{conj:dgf-edges}--\ref{conj:min-weight}): DGF
    outputs $O(n)$ edges and $O(\wt(\MST(P)))$ weight, and (more
    strongly) \emph{every} edge-minimal $t$-spanner in the plane has
    these same bounds. We support these conjectures experimentally and
    discuss why the weight conjecture for arbitrary minimal spanners
    is more delicate than the sparsity conjecture.
  \item \textbf{A faster implementation.} We show that the $\Theta(m)$
    APSP calls of the naive DGF can be reduced to $O(k\log m)$, where
    $m=\binom{n}{2}$ is the input size and $k$ is the output size
    (Algorithms~\ref{alg:dgf-bisect-main}--\ref{alg:bisect}). The
    technique is a recursive binary search that certifies contiguous
    batches of removable edges with a single APSP call.
  \item \textbf{Empirical comparison.} On random point sets in the
    unit square, with $n\in\{1000,\,2000,\,5000\}$ and
    $t\in\{1.001,\,1.01,\,1.1,\,1.5\}$, DGF and its variants produce
    spanners with edge counts that match the greedy spanner to
    within a fraction of a percent, but with total weight that is
    consistently $4$--$7\%$ lower, while preserving the exact stretch
    guarantee. We also examine the role of preprocessing: running DGF
    on top of a Yao graph leaves the output essentially unchanged, and
    yields a large runtime speedup for moderate $t$, but loses its
    advantage for very strict $t$ (because Yao itself is then no
    longer sparse), whereas the $\sqrt{t}$-Greedy preprocessor remains
    effective in both regimes.
\end{itemize}

\paragraph{Technical overview.}
The key observation behind DGF is that removing the \emph{longest}
redundant edge first allows shorter edges to inherit its role: a long
edge $(p,q)$ that can be detoured through intermediate points
$p\!\to\!x\!\to\!q$ at stretch $\le t$ is almost always detoured
through edges that are \emph{themselves} part of shorter redundant
detour paths, which DGF will still get to examine later. Greedy, by
contrast, commits to short edges first and can leave long edges
stranded. This ordering heuristic is so aggressive that in most
iterations a large prefix of edges is simultaneously removable; our
binary-search implementation exploits exactly this by bisecting on
prefixes, paying one APSP call per bisection rather than per edge.

For the general minimality question we treat the result of DGF as an
instance of a more general phenomenon: on any input, the set of
edge-minimal $t$-spanners forms a combinatorial family that is poorly
understood. We conjecture that this family always lies inside an
$O(n)$-edge, $O(\wt(\MST))$-weight envelope; a proof of either
conjecture would extend the existential-optimality results of Filtser
and Solomon from the single spanner produced by greedy to the entire
minimality class.

\paragraph{Paper organization.}
Section~\ref{sec:algorithm} defines DGF and analyses its time and space
complexity. Section~\ref{sec:minimal-def} defines minimal
$t$-spanners, and Section~\ref{sec:filtering} shows how DGF can be
used as a filter. Section~\ref{sec:conjectures} collects the four
sparsity and weight conjectures, for DGF and for arbitrary minimal
$t$-spanners.
Section~\ref{sec:binary-search} presents the recursive binary-search
implementation. Section~\ref{sec:experiments} reports experiments,
and Section~\ref{sec:related-work} surveys related work.
Section~\ref{sec:open-problems} concludes with open problems.
Auxiliary material---including the proof of
Theorem~\ref{thm:filter-minimal}, the per-regime experimental tables,
and an additional conjecture aimed at
Conjecture~\ref{conj:min-weight}---appears in the appendix.

\section{The Descending Greedy Filter Algorithm}
\label{sec:algorithm}

\begin{algorithm*}[h]
\caption{Descending Greedy Filter (DGF)}
\label{alg:dgf}
\begin{algorithmic}[1]
\Statex \textbf{Input:}
\Statex \hspace{1cm} $P$: a set of $n$ points in $\mathbb{R}^d$.
\Statex \hspace{1cm} $t$: a real constant with $t>1$.
\Statex \textbf{Output:}
\Statex \hspace{1cm} A set of edges $E$ such that $(P,E)$ is a $t$-spanner of $P$.

\State Sort the $\binom{n}{2}$ unordered pairs of distinct points in \textbf{non-ascending} order of
       their Euclidean distances (breaking ties arbitrarily) and store them in a list $L$.
\State $E \leftarrow L$
\For{each unordered pair $(p,q) \in L$ in \textbf{non-ascending} order}
    \State $E \leftarrow E \setminus \{(p,q)\}$
    \For{each unordered pair $(r,s) \in \binom{P}{2}$}
        \If{$\delta_{E}(r,s) > t \cdot |rs|$}
            \State $E \leftarrow E \cup \{(p,q)\}$
            \State \textbf{break}
        \EndIf
    \EndFor
\EndFor
\State \Return $E$
\end{algorithmic}
\end{algorithm*}

\subsection{Correctness}
At every point during the execution the current edge set $E$ is a
$t$-spanner of $P$: the algorithm starts with the complete graph (which
is trivially a $t$-spanner) and only removes an edge when the resulting
graph is still a $t$-spanner. Hence the returned $E$ is a $t$-spanner.

\subsection{Time and space complexity}
Let $m=\binom{n}{2}=\Theta(n^2)$ be the initial number of edges, and
let $T_{\APSP}(n,m)=O(n(m+n\log n))$ be the cost of a single
all-pairs shortest path computation (Dijkstra with a Fibonacci heap
from every source).

Sorting the $m$ edges costs $O(m\log m)=O(n^2 \log n)$.
For every edge we run one $\APSP$ to certify that its removal
preserves the stretch, giving total time
$O\bigl(m\cdot n(m+n\log n)\bigr) \;=\; O(n^5)$
when $m=\Theta(n^2)$, i.e.\ on the complete graph. The space required
to store $L$, $E$ and the APSP data structures is $O(m)=O(n^2)$.

\section{Minimal \texorpdfstring{$t$}{t}-spanners}
\label{sec:minimal-def}

\begin{definition}[minimal $t$-spanner]\label{def:minimal}
Let $P$ be a set of points and $t>1$. An edge set $E\subseteq \binom{P}{2}$
is a \emph{minimal $t$-spanner} of $P$ if
\begin{enumerate}
    \item $(P,E)$ is a $t$-spanner of $P$, and
    \item for every edge $(p,q)\in E$, the set $(E \setminus \{(p,q)\})$ is not a $t$-spanner for $P$.
\end{enumerate}
Equivalently:
\[
   \forall (p,q)\in E,\; \exists (r,s)\in P\times P\;:\;
   \delta_{E\setminus\{(p,q)\}}(r,s) > t\cdot|rs|.
\]
\end{definition}

Informally, a $t$-spanner is \emph{minimal} (also called \emph{edge-minimal})
when no edge can be removed without violating the stretch guarantee. The
greedy spanner is minimal by construction; most classical constructions
($\Theta$-graphs, Yao graphs, WSPD-based spanners) are not.

\section{Minimizing an arbitrary \texorpdfstring{$t$}{t}-spanner with DGF}
\label{sec:filtering}

The DGF procedure of Algorithm~\ref{alg:dgf} can be applied not only to
the complete graph but to any input edge set that already forms a
$t$-spanner. The only change is that the initial list $L$ contains
only the edges of the input $t$-spanner.

\begin{algorithm*}[h]
\caption{Descending Greedy Filter on a $t$-spanner}
\label{alg:dgf-filter}
\begin{algorithmic}[1]
\Statex \textbf{Input:}
\Statex \hspace{1cm} $P$: a set of $n$ points in $\mathbb{R}^d$.
\Statex \hspace{1cm} $t$: a real constant with $t>1$.
\Statex \hspace{1cm} $E_i$: an edge set such that $(P,E_i)$ is a $t$-spanner of $P$.
\Statex \textbf{Output:}
\Statex \hspace{1cm} A set of edges $E$ with $E\subseteq E_i$ such that $(P,E)$ is a minimal $t$-spanner of $P$.

\State Sort $E_i$ in \textbf{non-ascending} order by distance into a list $L$.
\State $E \leftarrow L$
\For{each unordered pair $(p,q) \in L$ in \textbf{non-ascending} order}
    \State $E \leftarrow E \setminus \{(p,q)\}$
    \For{each unordered pair $(r,s) \in \binom{P}{2}$}
        \If{$\delta_{E}(r,s) > t \cdot |rs|$}
            \State $E \leftarrow E \cup \{(p,q)\}$
            \State \textbf{break}
        \EndIf
    \EndFor
\EndFor
\State \Return $E$
\end{algorithmic}
\end{algorithm*}

\begin{theorem}\label{thm:filter-minimal}
Let $E_i$ be a $t$-spanner of $P$, and let $E$ be the output of
Algorithm~\ref{alg:dgf-filter} on input $(P,t,E_i)$. Then
(1) $E \subseteq E_i$, (2) $(P,E)$ is a $t$-spanner of $P$, and
(3) $E$ is a minimal $t$-spanner of $P$ in the sense of
Definition~\ref{def:minimal}.
\end{theorem}

A proof appears in Appendix~\ref{app:filter-proof}. The complexity
analysis is identical to that of Algorithm~\ref{alg:dgf}: with
$m_i=|E_i|$, sorting takes $O(m_i\log m_i)$ and the inner loop
$O(m_i\cdot n(m_i+n\log n))$ time. When $E_i$ is already a sparse
$t$-spanner ($m_i=O(n)$), this simplifies to $O(n^3 \log n)$.

\section{Sparsity and Weight Conjectures}
\label{sec:conjectures}

We now state our main conjectures. Throughout this section, $P$ is a
set of $n$ points in the plane and $t>1$ is a constant; the hidden
constants in the $O$-notation may depend on $t$. The conjectures come
in two flavors. The first pair concerns the specific output of DGF.
The second pair concerns \emph{any} minimal $t$-spanner in the sense
of Definition~\ref{def:minimal}.

\paragraph{Conjectures for DGF.}
Let $G_{\mathrm{DGF}}=(P,E_f)$ be the output of Algorithm~\ref{alg:dgf}
on $P$ with parameter $t$.

\begin{conj}[sparsity of DGF]\label{conj:dgf-edges}
$|E_f| \;=\; O(n)$.
\end{conj}

\begin{conj}[weight of DGF]\label{conj:dgf-weight}
$\wt(E_f) \;=\; O\!\bigl(\wt(\MST(P))\bigr)$.
\end{conj}

\paragraph{Conjectures for arbitrary minimal $t$-spanners.}
Let $E$ be any minimal $t$-spanner of $P$ in the sense of
Definition~\ref{def:minimal}, and assume $t < 2$.

\begin{conj}[sparsity of any minimal $t$-spanner]\label{conj:min-edges}
$|E| \;=\; O(n)$.
\end{conj}

\begin{conj}[weight of any minimal $t$-spanner]\label{conj:min-weight}
$\wt(E) \;=\; O\!\bigl(\wt(\MST(P))\bigr)$.
\end{conj}

Conjectures~\ref{conj:min-edges} and~\ref{conj:min-weight} strictly
generalize Conjectures~\ref{conj:dgf-edges} and~\ref{conj:dgf-weight}
respectively, since the output of DGF is minimal by
Theorem~\ref{thm:filter-minimal} (applied to $E_i=\binom{P}{2}$).
Likewise, the greedy spanner is minimal by construction, so the
conjectures are consistent with the classical bounds $O(n)$ edges
and $O(\wt(\MST(P)))$ weight that are known for the greedy spanner.
All four conjectures are supported by the experiments reported in
Section~\ref{sec:experiments}: in every tested instance the number of
edges of each minimal spanner we produce is close to $n$ and its
total weight is at most a small multiple of $\wt(\MST(P))$.
A further auxiliary conjecture, which we believe would substantially
advance a proof of Conjecture~\ref{conj:min-weight}, is stated in
Appendix~\ref{app:toward-min-weight}.

\section{Faster Implementation via Recursive Binary Search}
\label{sec:binary-search}

The na\"ive implementation of DGF (Algorithm~\ref{alg:dgf}) iterates
over all edges one by one, and for each edge performs an APSP
computation. This results in $\Theta(m)$ APSP calls, where
$m=\binom{n}{2}$. In practice, the vast majority of the long edges turn
out to be removable, and only a small subset is necessary. This
motivates a binary-search approach that can certify large batches of
removable edges with a single APSP call.

\paragraph{Key observation.}
Let $E$ be the current edge set (sorted in non-ascending order by
weight), and let $S \subseteq E$ be a contiguous batch of edges. If
removing all edges of $S$ from $E$ does not violate the $t$-spanner
condition---that is, if for every $r,s\in P$ we have
$\delta_{E\setminus S}(r,s) \le t\cdot|rs|$---then every edge in $S$
is removable. A single APSP call therefore certifies the removal of an
entire batch. When a violation is detected, we split the batch into two
halves and recurse on each. The pseudocode for the top-level driver
(Algorithm~\ref{alg:dgf-bisect-main}) and the recursive
\textsc{Bisect} routine (Algorithm~\ref{alg:bisect}) is given in
Appendix~\ref{app:bisect-pseudocode}.

\paragraph{Correctness and time complexity.}
The binary-search implementation produces the same output as
Algorithm~\ref{alg:dgf}: each edge is either certified as removable
(when the entire batch that contains it passes the violation check) or
identified as necessary (when the recursion reaches a single-edge base
case with a violation). Let $m$ denote the number of edges given to
\textsc{Bisect} at the top level, and let $k$ denote the number of
necessary edges in the output. Each necessary edge traces a
root-to-leaf path of depth at most $\lceil\log_2 m\rceil$ in the
recursion tree, so the total number of $\APSP$ calls is $O(k\log m)$.
With $T_{\APSP}(n,m)=O(n(m+n\log n))$ per call, the total time is
$O\!\bigl(m\log m + k\log m\cdot n(m+n\log n)\bigr)$. Since geometric
$t$-spanners are expected to satisfy $k=O(n)$ while $m=\Theta(n^2)$,
the speedup factor over the na\"ive implementation is
$\Omega(n/\log n)$, a near-linear improvement in the number of APSP
calls. On the complete graph this yields total time
$O(n^4 \log n)$; on a sparse $t$-spanner input ($m=O(n)$),
$O(n^3 \log^2 n)$.

\section{Experiments}
\label{sec:experiments}

\paragraph{Setup.}
All experiments are run on point sets drawn uniformly at random in the
unit square $[0,1]^2$, with fixed random $\mathtt{seed}=42$. We
compare \textbf{Greedy}, \textbf{$\sqrt{t}$-Greedy + DGF},
\textbf{Yao}, \textbf{Yao + DGF}, \textbf{$\sqrt{t}$-Yao +
$\sqrt{t}$-Greedy}, and \textbf{DGF}
(Algorithm~\ref{alg:dgf-bisect-main} starting from the complete
graph). For every instance we measure the number of edges, total
weight, weight ratio over the MST ($\wt(E)/\wt(\MST(P))$), max
and average vertex degree, the empirically measured maximum stretch,
and validity. Detailed per-regime tables (\ref{tab:n5000-t1.5}--\ref{tab:n1000-t1.001})
appear in Appendix~\ref{app:tables}.

\paragraph{Summary across regimes.}
To make the DGF-vs.-greedy comparison precise and uniform across the
four regimes, Table~\ref{tab:summary-percent} reports, for each
metric and each $(n,t)$ pair, the percentage change between the
\emph{best} DGF-based algorithm (over $\{\text{DGF},\,$Yao + DGF$,\,
\sqrt{t}\text{-Greedy} + \text{DGF}\}$) and the \emph{best}
greedy-based algorithm (over $\{$Greedy$,\,\sqrt{t}\text{-Yao} +
\sqrt{t}\text{-Greedy}\}$). Negative numbers favor the DGF family.
In our data the best greedy-based algorithm turns out to be plain
Greedy on every metric in every configuration.

\begin{table}[h]
    \centering
    \small
    \begin{tabular}{lrrrrr}
        \toprule
        Metric  & $t=1.5$ & $t=1.1$ & $t=1.01$ & $t=1.001$ & Avg. \\
        \midrule
        Edges   & $-0.76$ & $-0.44$ & $-0.29$ & $-0.28$ & $-0.44$ \\
        Weight  & $-6.91$ & $-5.86$ & $-4.40$ & $-3.98$ & $-5.28$ \\
        WtRatio & $-6.67$ & $-5.86$ & $-4.39$ & $-3.98$ & $-5.22$ \\
        MaxDeg  & $0.00$  & $0.00$  & $0.00$  & $+4.27$ & $+1.07$ \\
        AvgDeg  & $-0.53$ & $-0.45$ & $-0.29$ & $-0.27$ & $-0.39$ \\
        \bottomrule
    \end{tabular}
    \caption{Percentage change in each metric: best DGF-based output
    minus best greedy-based output, normalized by the best
    greedy-based output. Negative numbers favor the DGF family.}
    \label{tab:summary-percent}
\end{table}

\paragraph{Edges (and average degree).}
The best DGF-based algorithm in every cell is $\sqrt{t}$-Greedy +
DGF; it shaves between $0.28\%$ and $0.76\%$ off greedy's edge
count, averaging $0.44\%$. Average degree mirrors this finding. In
absolute terms these are all changes of well under one percent; for
practical purposes we view edge-minimal $t$-spanners of random
Euclidean inputs as having an edge count essentially determined by
the input, largely independent of which edge-minimization strategy
produces them. This is consistent with
Conjecture~\ref{conj:min-edges}.

\paragraph{Weight (and weight ratio).}
On total weight, the best DGF-based algorithm is consistently and
measurably lighter than greedy: the gap shrinks monotonically from
$-6.91\%$ at $t=1.5$ to $-3.98\%$ at $t=1.001$, averaging $-5.28\%$.
The weight ratio over the MST drops by essentially the same amount.
Descending processing, which removes the longest redundant edge
first, systematically leaves a slightly lighter edge-minimal
representative of the set of all edge-minimal $t$-spanners on the
input. This is the main empirical advantage of DGF over the
classical greedy spanner.

\paragraph{Maximum degree.}
The best DGF-based MaxDeg ties greedy exactly at $t=1.5,\;1.1$ and
$1.01$, and at $t=1.001$ exceeds greedy by $4.27\%$ ($122$ vs.\
$117$). The average MaxDeg increase is $+1.07\%$.

\paragraph{Yao preprocessing.}
The Yao graph is far from edge-minimal but provably contains a
$t$-spanner, so feeding it to DGF is safe. Empirically, the DGF
outputs from Yao are either identical to the full-input DGF output
(at $t=1.01$ and $t=1.001$) or differ by a handful of edges and a
negligible amount of weight: preprocessing is essentially neutral on
output quality. On runtime, for moderate $t$ ($t\ge 1.01$) Yao
prunes most candidate edges and the binary-search DGF gets a large
practical speedup; at $t=1.001$, Yao's cones become so narrow that
it retains $93\%$ of $\binom{n}{2}$ edges and the speedup vanishes.
The $\sqrt{t}$-Greedy preprocessor produces an $O(n)$-edge
graph and remains effective in both regimes; at $t=1.001$ it gives
roughly an order-of-magnitude speedup over pure DGF.

\paragraph{Support for the conjectures.}
The weight ratios over the MST stay bounded and scale with $t$ as
classical theory predicts: $\rho\approx 2.8$ for $t=1.5$,
$\rho\approx 10$ for $t=1.1$, $\rho\approx 57$ for $t=1.01$ and
$\rho\approx 267$ for $t=1.001$, consistent with the $1/(t-1)$
behavior observed for the greedy spanner. Edge counts remain linear
in $n$ across all four $(n,t)$ pairs. Together, these observations
give direct empirical support to
Conjectures~\ref{conj:dgf-edges}--\ref{conj:min-weight}: every
edge-minimal $t$-spanner we compute has a linear number of edges
and weight $O(\wt(\MST))$.

\section{Related Work}
\label{sec:related-work}

DGF sits at the intersection of two lines of work on geometric
spanners: (i) constructive algorithms that build small or light
$t$-spanners directly (greedy, Yao, $\Theta$-graphs, WSPD-based), and
(ii) structural properties characterizing when a $t$-spanner is sparse
and/or light (the $t$-leapfrog property, existential optimality of
greedy). A comprehensive reference is the monograph of Narasimhan and
Smid~\cite{NS07}.

\paragraph{Greedy and the leapfrog property.}
The greedy spanner was introduced by Alth\"ofer et
al.~\cite{ADDJS93}. DGF is the natural reverse: instead of starting
from the empty graph and inserting short edges first, DGF starts from
a $t$-spanner and removes long edges first. Both algorithms terminate
at an edge-minimal $t$-spanner, but on the same input they produce
different subgraphs in general. Sparsity and weight bounds for greedy
in Euclidean space are due to Chandra et al.~\cite{CDNS95}; the
$t$-leapfrog property and its weight bound were introduced by Das et
al.~\cite{DHN93,DNS95}. Filtser and Solomon~\cite{FS20} proved that
greedy is existentially (near-)optimal. DGF is interesting precisely
because its outputs are edge-minimal but do \emph{not} automatically
satisfy the leapfrog property; whether they are nevertheless light is
the subject of Conjectures~\ref{conj:dgf-weight}
and~\ref{conj:min-weight}.

\paragraph{Cone-based and Delaunay spanners.}
The cone-based family begins with Yao~\cite{Yao82}; the
$\Theta$-graph variant is due to Clarkson~\cite{Clarkson87} and
Keil~\cite{Keil88}, with $t$-spanner analyses by Keil and
Gutwin~\cite{KG92}. The Delaunay triangulation was shown to be a
$t$-spanner by Dobkin et al.~\cite{DFS90}. These deterministic,
oblivious constructions are usually \emph{not} edge-minimal, which is
why running DGF on top of a Yao graph drops edges aggressively in our
experiments.

\paragraph{WSPD and light spanners.}
The Well-Separated Pair Decomposition was introduced by Callahan and
Kosaraju~\cite{CK95}. Light planar $t$-spanners were given by
Levcopoulos and Lingas~\cite{LL92}; light $(1+\varepsilon)$-spanners
in arbitrary fixed dimension by Das, Narasimhan and Salowe~\cite{DNS95},
with a near-quadratic implementation in Das and Narasimhan~\cite{DN97}.
Borradaile, Le and Wulff-Nilsen~\cite{BLW19} extended the optimality
of greedy to doubling metrics, and Le and Solomon~\cite{LS25} proved
tight bounds on the size and lightness of Euclidean
$(1+\varepsilon)$-spanners.

\paragraph{Faster greedy and approximate-greedy spanners.}
The exact greedy spanner has been sped up to $O(n^2\log n)$ time by
Bose et al.~\cite{BCFMS10} and to $O(n)$ space with $O(n^2\log^2 n)$
time by Alewijnse et al.~\cite{ABBB15}; an $O(n\log n)$
\emph{approximate}-greedy variant is due to Gudmundsson et
al.~\cite{GLN02}. Our recursive binary-search variant of DGF
(Section~\ref{sec:binary-search}) is in the same engineering spirit:
it batches edges and uses a logarithmic number of APSP rounds to
certify many edges at once.

\paragraph{Edge-minimality, filtering, and empirical comparisons.}
We are not aware of prior work proposing the descending-order filter
algorithm we study. DGF can be viewed as the natural spanner-analogue
of the reverse-delete algorithm for spanning trees, with ``connected''
replaced by ``$t$-spanner of the input''. Computing the smallest
$t$-spanner is NP-hard~\cite{Cai94} and inapproximable within
$(\log n)^{1-\delta}$~\cite{Kortsarz01}; an ILP-based exact algorithm
for the minimum-weight $t$-spanner is given by Sigurd and
Zachariasen~\cite{SZ04}. Edge-minimality is a strictly weaker,
polynomial-time attainable objective. On the empirical side, Farshi
and Gudmundsson~\cite{FG09} carried out the first systematic
comparison of standard constructions, and Alewijnse et
al.~\cite{ABBB15} extended these experiments to point sets of up to
$10^6$ vertices.

\section{Conclusion and Open Problems}
\label{sec:open-problems}

We introduced the Descending Greedy Filter, a construction dual to the
classical greedy spanner, together with a general filtering procedure
that turns any $t$-spanner into an edge-minimal one. We gave a
recursive binary-search implementation that reduces the number of
$\APSP$ calls from $\Theta(m)$ to $O(k\log m)$. Empirically, on random
point sets with $n\in\{1000,\,2000,\,5000\}$ and
$t\in\{1.001,\,1.01,\,1.1,\,1.5\}$, DGF and its variants produce
spanners whose edge counts match the greedy spanner to within a
fraction of a percent while using $4$--$7\%$ less total weight. The
choice of preprocessing affects only runtime and not output quality:
running DGF on top of a Yao graph is a large practical speedup for
moderate $t$ but loses its advantage at very strict $t$ (because Yao
itself ceases to be sparse), while the $\sqrt{t}$-Greedy preprocessor
remains effective in both regimes.

The main questions left open by this work are the four conjectures of
Section~\ref{sec:conjectures}. In order of decreasing confidence:
(1) Sparsity of DGF (Conjecture~\ref{conj:dgf-edges}); our experiments
show the number of edges is consistently below $2n$ for the tested
parameters. (2) Sparsity of every minimal $t$-spanner in the plane
(Conjecture~\ref{conj:min-edges}), which would extend Filtser and
Solomon's~\cite{FS20} existential-optimality result from the greedy
spanner to the entire class of edge-minimal spanners. (3) Weight of
DGF (Conjecture~\ref{conj:dgf-weight}); descending-order processing
appears to enforce a leapfrog-like property on the output, but we do
not have a proof. (4) Weight of every minimal $t$-spanner
(Conjecture~\ref{conj:min-weight}); this is the weakest and may in
fact be false, since the classical weight bound for
greedy~\cite{CDNS95,DHN93,DNS95} rests on the $t$-leapfrog property,
which is not implied by edge-minimality alone.

A further direction is to replace the $O(k\log m)$ $\APSP$ calls of
the binary-search implementation by a genuinely near-linear-time
algorithm (analogous to~\cite{GLN02,BCFMS10,ABBB15} in the greedy
setting) that incrementally maintains the shortest-path structure as
edges are removed. Finally, it would be interesting to combine DGF
with a non-trivial \emph{starting} spanner (a $\Theta$-graph or a WSPD
spanner) and study the distribution of edge-minimal outputs obtained
by different starting points.


%---------------------------- Bibliography -------------------------------

% Per CCCG style: when ready for final submission, replace the
% \bibliography call with the contents of the .bbl file generated by
% bibtex (using \begin{thebibliography}{99} ... \end{thebibliography}),
% sorted alphabetically. For drafting we keep \bibliography{references}.
\small
\bibliographystyle{abbrv}
\bibliography{references}


\newpage
\appendix

\section{Proof of Theorem~\ref{thm:filter-minimal}}
\label{app:filter-proof}

Property (1) is immediate, since the algorithm only removes edges.

Property (2) follows from the invariant that $E$ is a $t$-spanner at
every iteration: the algorithm starts with the $t$-spanner $E_i$, and
whenever an edge is removed we verify that the result is still a
$t$-spanner; otherwise the edge is restored.

Property (3). Let $(p,q)\in E$ be an edge in the output. We must
exhibit a pair $(r,s)\in P\times P$ such that
$\delta_{E\setminus\{(p,q)\}}(r,s)>t\cdot|rs|$.

Denote by $E'$ the state of $E$ at the moment the main loop inspected
the edge $(p,q)$. Because the algorithm only removes edges after the
inspection of $(p,q)$, we have $E\subseteq E'$.

The fact that $(p,q)\in E$ means that the inner violation check
succeeded, i.e.\ there exist $(r,s)\in P\times P$ with
$\delta_{E'\setminus\{(p,q)\}}(r,s) > t\cdot|rs|$.
Since $E\subseteq E'$, removing more edges can only increase shortest
path distances, hence
$\delta_{E\setminus\{(p,q)\}}(r,s) \ge
\delta_{E'\setminus\{(p,q)\}}(r,s) > t\cdot|rs|$,
which is exactly the minimality condition for $(p,q)$.
\hfill$\QED$


\section{A Conjecture toward Conjecture~\ref{conj:min-weight}}
\label{app:toward-min-weight}

We state an additional conjecture which, if established, we believe
would imply (or substantially advance a proof of)
Conjecture~\ref{conj:min-weight}.

\paragraph{Setup.}
Let $E_f$ be the output of the DGF algorithm on the complete graph on
$P$. Let $E = \{(p,q) \in E_f : \delta_2(p,q) \leq t \cdot |pq|\}$,
i.e.\ all the edges that fail the $t$-leapfrog property.

\paragraph{Notation.}
For each edge $(p,q) \in E$, denote
\begin{align*}
  C_{pq} &\;=\; \Bigl\{(a,b) \in E \;\Bigm|\;
    \delta_{E_f\setminus(p,q)}(a,b) > t \cdot |ab|\Bigr\}, \\
  (a,b)_{pq} &\;=\; \arg\min_{(a,b)\,\in\, C_{pq}} |ab|, \\
  T_{pq} &\;=\; \delta_{E_f}(a,p) + |pq| + \delta_{E_f}(q,b),
\end{align*}
where $(a,b)=(a,b)_{pq}$.

\paragraph{Construction of $F$.}
Initialize $F \gets \emptyset$ and $E^* \gets E$. Sort $E^*$ in
non-ascending order. For every $(p,q) \in E^*$:
$F \gets F \cup T_{pq}$, $E^* \gets E^* \setminus T_{pq}$.

\begin{conj}[auxiliary conjecture toward weight of any minimal $t$-spanner]
\label{conj:toward-min-weight}
With the construction above, $\wt(F) = \wt(\MST(P))$.
\end{conj}

\paragraph{Why this would help.}
If Conjecture~\ref{conj:toward-min-weight} holds, then every minimal
$t$-spanner $E$ admits a witness collection $\{T_{pq}\}_{(p,q)\in E}$
whose total weight is at most $\wt(\MST(P))$; combined with the fact
that each edge $(p,q)\in E$ contributes its own length $|pq|$ to some
$T_{pq}$, this would yield $\wt(E) = O(\wt(\MST(P)))$ and hence
Conjecture~\ref{conj:min-weight}.


\section{Pseudocode of the recursive binary-search implementation}
\label{app:bisect-pseudocode}

This appendix gives the pseudocode for the binary-search variant of
DGF described in Section~\ref{sec:binary-search}.
Algorithm~\ref{alg:dgf-bisect-main} is the top-level driver, and
Algorithm~\ref{alg:bisect} is the recursive routine that confirms or
refines a contiguous batch of candidate-removable edges with a single
$\APSP$ call.

\begin{algorithm}[h]
\caption{DGF with recursive binary search}
\label{alg:dgf-bisect-main}
\begin{algorithmic}[1]
\Statex \textbf{Input:} $P$, $t$ as in Algorithm~\ref{alg:dgf}.
\Statex \textbf{Output:} a $t$-spanner $E$ of $P$ identical to the output of Algorithm~\ref{alg:dgf}.

\State Sort the $\binom{n}{2}$ pairs in \textbf{non-ascending} order by distance into a list $L$.
\State $E \leftarrow L$
\State $E \leftarrow \textsc{Bisect}(\emptyset,\, L,\, P,\, t)$
\State \Return $E$
\end{algorithmic}
\end{algorithm}

\begin{algorithm*}[h]
\caption{\textsc{Bisect}$(E,\,S,\,P,\,t)$}
\label{alg:bisect}
\begin{algorithmic}[1]
\Statex \textbf{Input:}
\Statex \hspace{1cm} $E$: current edge set (i.e.\ $S$ has been removed from $E$ before the call).
\Statex \hspace{1cm} $S$: batch of candidate edges to confirm as removable, sorted in non-ascending order by weight.
\Statex \hspace{1cm} $P$: point set.
\Statex \hspace{1cm} $t$: stretch parameter.
\Statex \textbf{Output:} updated $E$ in which the removable edges of $S$ are discarded and the necessary edges are restored.

\If{$S = \emptyset$}
    \State \Return $E$
\EndIf
\State Compute APSP on $(P,E)$.
\If{$\delta_E(r,s) \le t \cdot |rs|$ for every $(r,s)\in P\times P$}
    \State \Return $E$ \Comment{No violation: every edge in $S$ is removable}
\EndIf
\If{$|S|=1$}
    \State $E \leftarrow E \cup S$ \Comment{Single necessary edge: restore it}
    \State \Return $E$
\EndIf
\State Split $S$ into $S_1$ (first, longer half) and $S_2$ (second, shorter half).
\State $E \leftarrow E \cup S$ \Comment{Restore all edges in $S$}
\State $E \leftarrow E \setminus S_1$
\State $E \leftarrow \textsc{Bisect}(E,\, S_1,\, P,\, t)$
\State $E \leftarrow E \setminus S_2$
\State $E \leftarrow \textsc{Bisect}(E,\, S_2,\, P,\, t)$
\State \Return $E$
\end{algorithmic}
\end{algorithm*}


\section{Detailed experimental tables}
\label{app:tables}

\begin{table*}[h]
    \centering
    \small
    \begin{tabular}{lrrrrrrr}
        \toprule
        Algorithm & Edges & Weight & WtRatio & MaxDeg & AvgDeg & MaxStrch & Valid \\
        \midrule
        Yao + DGF                             & 9\,377           & \textbf{128.07}  & \textbf{2.80} & \textbf{7} & 3.75          & 1.5000           & \checkmark \\
        $\sqrt{t}$-Greedy + DGF               & \textbf{9\,314}  & 128.91           & 2.81          & \textbf{7} & \textbf{3.73} & 1.4999           & \checkmark \\
        Greedy                                & 9\,385           & 137.58           & 3.00          & \textbf{7} & 3.75          & 1.4999           & \checkmark \\
        $\sqrt{t}$-Yao + $\sqrt{t}$-Greedy    & 14\,269          & 255.68           & 5.58          & 11         & 5.71          & 1.2754           & \checkmark \\
        Yao                                   & 66\,839          & 2\,302.90        & 50.26         & 52         & 26.74         & \textbf{1.1647}  & \checkmark \\
        \bottomrule
    \end{tabular}
    \caption{Summary for $n=5000$, $t=1.5$, $\mathtt{seed}=42$.}
    \label{tab:n5000-t1.5}
\end{table*}

\begin{table*}[h]
    \centering
    \small
    \begin{tabular}{lrrrrrrr}
        \toprule
        Algorithm & Edges & Weight & WtRatio & MaxDeg & AvgDeg & MaxStrch & Valid \\
        \midrule
        Yao + DGF                             & 22\,191           & \textbf{463.89}   & \textbf{10.12} & 17           & 8.88          & 1.1000          & \checkmark \\
        $\sqrt{t}$-Greedy + DGF               & \textbf{22\,127}  & 467.67            & 10.21          & \textbf{16}  & \textbf{8.85} & 1.1000          & \checkmark \\
        Greedy                                & 22\,225           & 492.74            & 10.75          & \textbf{16}  & 8.89          & 1.1000          & \checkmark \\
        $\sqrt{t}$-Yao + $\sqrt{t}$-Greedy    & 32\,320           & 858.65            & 18.74          & 25           & 12.93         & 1.0488          & \checkmark \\
        Yao                                   & 224\,442          & 14\,250.23        & 311.00         & 137          & 89.78         & \textbf{1.0378} & \checkmark \\
        \bottomrule
    \end{tabular}
    \caption{Summary for $n=5000$, $t=1.1$, $\mathtt{seed}=42$.}
    \label{tab:n5000-t1.1}
\end{table*}

\begin{table*}[h]
    \centering
    \small
    \begin{tabular}{lrrrrrrr}
        \toprule
        Algorithm & Edges & Weight & WtRatio & MaxDeg & AvgDeg & MaxStrch & Valid \\
        \midrule
        DGF                                   & 27\,714           & \textbf{1\,647.42} & \textbf{56.85}  & \textbf{44}  & 27.71          & 1.0100          & \checkmark \\
        Yao + DGF                             & 27\,714           & \textbf{1\,647.42} & \textbf{56.85}  & \textbf{44}  & 27.71          & 1.0100          & \checkmark \\
        $\sqrt{t}$-Greedy + DGF               & \textbf{27\,659}  & 1\,657.07          & 57.18           & \textbf{44}  & \textbf{27.66} & 1.0100          & \checkmark \\
        Greedy                                & 27\,739           & 1\,723.16          & 59.46           & \textbf{44}  & 27.74          & 1.0100          & \checkmark \\
        $\sqrt{t}$-Yao + $\sqrt{t}$-Greedy    & 38\,837           & 2\,864.12          & 98.83           & 57           & 38.84          & 1.0050          & \checkmark \\
        Yao                                   & 571\,029          & 156\,774.77        & 5\,409.63       & 860          & 571.03         & \textbf{1.0008} & \checkmark \\
        \bottomrule
    \end{tabular}
    \caption{Summary for $n=2000$, $t=1.01$, $\mathtt{seed}=42$.}
    \label{tab:n2000-t1.01}
\end{table*}

\begin{table*}[h]
    \centering
    \small
    \begin{tabular}{lrrrrrrr}
        \toprule
        Algorithm & Edges & Weight & WtRatio & MaxDeg & AvgDeg & MaxStrch & Valid \\
        \midrule
        DGF                                   & 38\,633          & \textbf{5\,556.17} & \textbf{267.16}  & 127           & 77.27          & 1.0010          & \checkmark \\
        Yao + DGF                             & 38\,633          & \textbf{5\,556.17} & \textbf{267.16}  & 127           & 77.27          & 1.0010          & \checkmark \\
        $\sqrt{t}$-Greedy + DGF               & \textbf{38\,574} & 5\,588.94          & 268.74           & 122           & \textbf{77.15} & 1.0010          & \checkmark \\
        Greedy                                & 38\,682          & 5\,786.30          & 278.23           & \textbf{117}  & 77.36          & 1.0010          & \checkmark \\
        $\sqrt{t}$-Yao + $\sqrt{t}$-Greedy    & 52\,375          & 9\,167.11          & 440.79           & 164           & 104.75         & 1.0005          & \checkmark \\
        Yao                                   & 465\,760         & 235\,769.69        & 11\,336.68       & 982           & 931.52         & \textbf{1.0001} & \checkmark \\
        \bottomrule
    \end{tabular}
    \caption{Summary for $n=1000$, $t=1.001$, $\mathtt{seed}=42$.}
    \label{tab:n1000-t1.001}
\end{table*}

\end{document}
